% !TEX TS-program = lualatex

\documentclass[border = 3pt]{standalone}

%%%
% Des couleurs faciles d'emploi via le package \xcolor!
%%%
\usepackage[svgnames]{xcolor}

%%%
% La bibliothèque \luadraw allie une facilité d’utilisation
% à un rendu particulièrement soigné.
%%%
\usepackage[3d]{luadraw}

\begin{document}

\begin{luadraw*}{name = orientation-n-box-axes-full}
local ld = luadraw

local M = ld.pt3d.M

------
-- Nos ingrédients de base.
------
local pi   = math.pi
local sin  = math.sin
local cos  = math.cos
local tanh = math.tanh

------
-- ¨Def de la zone graphique.
------
local graphview = ld.graph3d:new{
  window3d = {-3, 6, -4, 6, -2, 2},
  viewdir  = {210, 70},  -- Un point de vue éclairant.
  adjust2d = true,
  size     = {14, 14},
  bbox     = false
}

------
-- Ajout d'un repère de type "boîte ouverte" avant la surface
-- (ne jamais oublier que ¨luadraw n'est pas un moteur de rendu
-- ¨3d).
------
graphview:Dboxaxes3d({
  grid      = true,
  gridcolor = "gray",
  fillcolor = "lightgray"
})

------
-- ¨Def et tracé de notre ¨surf ¨param en mode monocolore.
------
local facets = ld.surface(
  function(u, v)
    return M(
      2*(1 + cos(v))*cos(u),
      2*(1 + cos(v))*sin(u),
      2*sin(v)*tanh(u - pi)
    )
  end,
  0, 2*pi, 0, 2*pi,
  {51, 15}
)

graphview:Dfacet(
  facets,
  {
    mode  = ld.mShadedOnly,
    color = "ForestGreen"
  }
)

------
-- Montrons le résultat de notre oeuvre.
------
graphview:Show()
\end{luadraw*}

\end{document}
